% !TEX program = lualatex
\documentclass[10pt, leqno]{ltjsarticle}
% \documentclass[10pt]{ltjsarticle} % 日本語をそのまま扱えるクラス
% \usepackage[margin=28mm]{geometry}
\usepackage{amsmath,amssymb,mathtools}
\usepackage{amsthm}
\usepackage{tikz}
\usetikzlibrary{arrows.meta,positioning,calc}

\usepackage[
  top=15mm,     % 上余白
  bottom=15mm,  % 下余白
  left=10mm,    % 左余白
  right=10mm,    % 右余白
  landscape,
  margin=10mm
]{geometry}
% \usepackage[landscape,margin=20mm]{geometry} % 余白はお好みで
\usepackage{paracol}
\columnratio{0.5,0.5}             % 左右を 1:1 に
\setlength{\columnseprule}{0.4pt} % 真ん中の縦線

% ---- 記号・マクロ ----
\newcommand{\xto}{\mathrel{\xrightarrow{}}}
\newcommand{\into}{\hookrightarrow}
\newcommand{\epi}{\twoheadrightarrow}
\newcommand{\pto}{\rightrightarrows}
\newcommand{\cls}[1]{\bigl(#1\bigr)}
\newcommand{\redjoin}{\wedge}
\newcommand{\shrink}{\phi}
\newcommand{\Ex}{E}        % suspension functor E
\newcommand{\Ox}{\Omega}   % loop functor Omega
\newcommand{\id}{\mathrm{id}}
% --- Toda bracket helpers ---
\newcommand{\toda}[1]{\{\,#1\,\}}              % {a,b,c}
\newcommand{\todai}[2][]{\{\,#2\,\}_{#1}}      % {a,b,c}_i   (#1 は添字, 空なら無し) 
\newcommand{\todas}[1]{\langle\,#1\,\rangle}
\newcommand{\gen}[1]{\{\,#1\,\}}       % ⟨generator⟩ 表記を { } で統一
\newcommand{\Ker}{\operatorname{Ker}}
\newcommand{\im}{\operatorname{Im}}
\newcommand{\hs}[2]{\bigl[#1,#2\bigr]}   % 一般のホモトピー集合
\newcommand{\pis}[2]{\pi_{#1}(#2)}       % 左が球面のときのホモトピー群
\newcommand{\pit}[2]{\pi_{#1}^{#2}}
\newcommand{\ext}[1]{\overline{#1}}     % extension
\newcommand{\coext}[1]{\widetilde{#1}}  % coextension
\newcommand{\z}{\mathbb{Z}}
\newcommand{\sq}[1]{Sq^{#1}}
\newcommand{\map}{\operatorname{Map}}

% ---- theorem-like environments ----
\numberwithin{equation}{section}
\theoremstyle{plain}
\newtheorem{theorem}{Theorem}[section]
\newtheorem{lemma}[theorem]{Lemma}
\newtheorem{proposition}[theorem]{Proposition}
\newtheorem{corollary}[theorem]{Corollary}
\theoremstyle{definition}
\newtheorem{definition}[theorem]{Definition}
\newtheorem{example}[theorem]{Example}
\newtheorem{notation}[theorem]{Notation}
\newtheorem{assumption}[theorem]{Assumption}
\theoremstyle{remark}
\newtheorem{remark}[]{Remark}

% \title{PART 1 \quad CHAPTER I\\[2mm] General Composition and Secondary Composition}
% \author{}
% \date{}

\begin{document}

% ========== CHAPTER IV ==========
\setcounter{section}{3}
\section{Suspension Sequence mod.\ 2}

\begin{paracol}{2}

よく知られている事実として, 
\begin{equation}\label{4.1}
\pis{i}{S^1}=0\ (i>1),\qquad
\pis{i}{S^n}=0\ (i<n),\qquad
\pis{n}{S^n}\cong \mathbb{Z}\gen{\iota_n} .
\end{equation}

また, 次のことも知られている (Serre \cite{Serre1951}) :
\begin{equation}\label{4.2}
% (4.2).\
\pis{i}{S^n}\ \text{は,} \ i=n \ \text{または} \ 
i=2n-1 \ \text{のとき以外は, 有限群である. }
\end{equation}

\(\pis{i}{S^n}\) の $p$-成分を, $\pis{i}{S^n;p}$ と書き, 
部分群 $\pi_i^n\subset\pis{i}{S^n}$ を, 次で定義する. 
\begin{equation}\label{4.3}
\pi_i^n :=
\begin{cases}
\pis{n}{S^n}, & i=n,\\[2pt]
\Ex^{-1}\bigl(\pis{2n}{S^{n+1};2}\bigr), & i=2n-1,\\[2pt]
\pis{i}{S^n;2}, & i\ne n,\,2n-1 .
\end{cases}
\end{equation}

群 $\pi^n_{2n-1}$ は, 次で具体的に与えられる. 
\begin{lemma}\label{lem:4.1}
% Lem4.1.\
$n$ が奇数のとき,  $\pi_{2n-1}^n=\pis{2n-1}{S^n;2} .$\\
$n$ が偶数で, $[\,\iota_{n-1},\iota_{n-1}\,]\ne0$ のとき, 
\(
\pi_{2n-1}^n\ \cong\ \mathbb{Z}\gen{\Delta(\iota_{2n+1})}\ \oplus \ 
\pis{2n-1}{S^n;2} .
\)\\
$n$ が偶数で, $[\,\iota_{n-1},\iota_{n-1}\,]=0$ のとき, 
\(
\pi_{2n-1}^n\ \cong\ \mathbb{Z}\gen{\alpha}\ \oplus \
\pis{2n-1}{S^n;2} .
\)\\
ここで, \(\alpha\) は 
\(H(\alpha)=\iota_{2n-1},\ E\alpha\in\pi_{2n}^{n+1}=\pi_{2n}(S^{n+1};2)\)
を満たす元である. \\
いずれの場合も, $[\,\iota_n,\iota_n\,]\in \pi_{2n-1}^n$ である. 
\end{lemma}

\noindent
\textbf{Proof.}\
まず, $n$ を奇数とする. \\
Whitehead product の anti-commutativity より,
$2\Delta(\iota_{2n+1})=2[\,\iota_n,\iota_n\,]=0 .$\\
(2.11) の完全性から,
\(
\Ker\bigl(\Ex:\pis{2n-1}{S^n}\to\pis{2n}{S^{n+1}}\bigr)
\subset\pis{2n-1}{S^n;2}.
\)\\
(4.2) より, $\pis{2n-1}{S^n}$ は有限群なので, 
\(
\pi_{2n-1}^n=\Ex^{-1}\bigl(\pis{2n}{S^{n+1};2}\bigr)=\pis{2n-1}{S^n;2} .
\)

% --- CHAPTER IV p.35 ---

次に, $n$ を偶数とし, (2.11) の完全列を考える. 
\[
\pis{2n-2}{S^{n-1}} \xrightarrow{\ \Ex\ } \pis{2n-1}{S^{n}}
\xrightarrow{\ H\ } \pis{2n-1}{S^{2n-1}}
\xrightarrow{\ \Delta\ } \pis{2n-3}{S^{n-1}}.
\]
$\Delta$ の象は $[\,\iota_{n-1},\iota_{n-1}\,]$ によって生成され, 
また, $2[\,\iota_{n-1},\iota_{n-1}\,]=0$ であるから, 
\[
\pis{2n-1}{S^{n}} \cong
\mathbb{Z}\gen{\beta}\ \oplus\ 
E\pis{2n-2}{S^{n-1}} .
\]
ここで,
\(E(\pis{2n-2}{n-1})\) は有限であり, \\
$[\,\iota_{n-1},\iota_{n-1}\,]=0$ の場合, $H(\beta)=\iota_{2n-1}$,\\
$[\,\iota_{n-1},\iota_{n-1}\,]\ne0$ の場合, $H(\beta)=2\iota_{2n-1}$ である. 

\switchcolumn

\noindent
\textbf{\cite{Serre1951}
J.-P.~Serre,}
\newblock Homologie singulière des espaces fibrés,
\newblock Ann. of Math., {\bf 54} (1951), 425--505.

\noindent
\([\beta, \alpha] = (-1)^{pq}[\alpha, \beta],\ 
\alpha\in\pis{p}{X},\ \beta\in\pis{q}{X}\).\\
\(n\) が奇数のとき,
\([\iota_n, \iota_n] = -[\iota_n, \iota_n]\)
なので, \(2[\iota_n, \iota_n] = 0.\)

\noindent
\textbf{(2.11).}\ 
\(\cdots\ \longrightarrow\ 
\pi_i(S^m)\xrightarrow{\,E\,}\pi_{i+1}(S^{m+1})\xrightarrow{\,H\,}
\pi_{i+1}(S^{2m+1})\xrightarrow{\,\Delta\,}\pi_{i-1}(S^m)\xrightarrow{\,E\,}\cdots
\)\\
これは $m$ が奇数, あるいは $i<3m-1$ のとき完全列である. \\
$m$ が偶数で, $i\ge 3m-1$ の場合には, (2.11) を
2-成分に関する完全列として考える. \\
\(\pis{2n+1}{S^{2n+1}}\xlongrightarrow{\Delta}
\pis{2n-1}{S^{n}}\xlongrightarrow{E}
\pis{2n}{S^{n+1}}\).\\
\(\ker E = \im \Delta = \gen{\Delta(\iota_{2n+1})}
\subset\pis{2n-1}{S^n; 2}\).

\noindent
\(2\)-成分に写る元は,
\(2\)-成分の元と, 無限位数の元と, kernel の元である.



\end{paracol}

\begin{paracol}{2}

\noindent
Prop2.7 により, 後者の場合には,
$\beta=\pm[\,\iota_n,\iota_n\,]$ が取れる. \\
さらに, $\Ker \Ex:\pis{2n-1}{S^{n}}\to\pis{2n}{S^{n+1}}$ は,
\(\pm[\iota_n, \iota_n]=\Delta(\iota_{2n+1})\) で生成されるので, 
\[
\begin{align*}
  \pit{2n-1}{n} = E^{-1}(\pis{2n}{S^{n+1}; 2})
  &= \Delta(\pis{2n+1}{S^{2n+1}}) \oplus (E \pis{2n-2}{n-1}; 2) \\
  &= \Delta(\pis{2n+1}{S^{2n+1}}) \oplus \pis{2n-1}{S^n;2} .
\end{align*}
\]
以上で, $n$ が偶数 かつ, 
$[\,\iota_{n-1},\iota_{n-1}\,]\ne0$ の場合のLem~\ref{lem:4.1} が示された. 

次に, $n$ 偶数で,
$[\,\iota_{n-1},\iota_{n-1}\,]=0$ の場合を考える. \\
\eqref{4.2} より, $\pis{2n}{S^{n+1}}$ は有限なので, 
ある整数 $t$ が存在して,
\((2t+1)E\beta\in\pis{2n}{S^{n+1};2}\). \\
ここで, 
\(
\alpha\ :=\ \beta- t \bigl(\pm[\,\iota_n,\iota_n\,]\bigr)
\)
と置く. 
符号は, Prop2.7 の $H[\,\iota_n,\iota_n\,]=\pm2\iota_{2n-1}$ による. \\
このとき, 
$E\alpha\in \pis{2n}{S^{n+1};2}$ かつ $H(\alpha)=\iota_{2n-1}$ となる. \\
上記の列の完全性により, 
\(
\pis{2n-1}{S^{n}}\ \cong\
\mathbb{Z}\gen{\alpha} \oplus E\pis{2n-2}{S^{n-1}} .
\)\\
$\Ex(\pis{2n-1}{S^n;2} \oplus \gen{\alpha}) = \pis{2n}{S^{n+1};2}$ 
であることも容易に分かる. 

差 \([\iota_n, \iota_n]-2\alpha\) について考える. \\
\(H([\iota_n, \iota_n] - 2\alpha) = 0\) なので, 
\([\iota_n, \iota_n] - 2\alpha \in E \pis{2n-2}{S^n-1}\) である. \\
\(E([\iota_n, \iota_n] - 2\alpha) = -2E\alpha \in \pis{2n}{S^{n+1}; 2}\)
であり, \\
\(E:\pis{2n-1}{S^n}\to\pis{2n}{S^{n+1}}\)
は, 
\(E\pis{2n-2}{S^{n-1}}\)
を 
\(\pis{2n}{S^{n+1}}\)
に同型に写すので, 
\[
[\iota_n, \iota_n] \in E\pis{2n-2}{S^n; 2} \oplus \gen{\alpha}
= \pis{2n-1}{S^n; 2} \oplus \gen{\alpha} .
\]
以上より, 
\(
\pit{2n-1}{n} = E\pis{2n}{S^{n+1}; 2}
= \pis{2n-1}{S^n; 2} \oplus \gen{\alpha} .
\)
\qed

次の完全列は, \(\pis{n+k}{S^n}\) の $2$-成分を計算するための主道具である. 
\begin{proposition}[]\label{prop:4.2}
% Prop4.2.\
次の列は完全列である：
\begin{equation}\label{4.4}
% (4.4)
\cdots \longrightarrow
\pit{i}{n}\ \xrightarrow{\ \Ex\ }\ \pit{i+1}{n+1}\
\xrightarrow{\ H\ }\ \pit{i+1}{2n+1}\
\xrightarrow{\ \Delta\ }\ \pit{i-1}{n}\
\xrightarrow{\ \Ex\ }\ \pit{i}{n+1}\ \xrightarrow{\ H\ }\ \cdots.
\end{equation}
\end{proposition}

\noindent
\textbf{Proof.}\
Th2.4 と (2.11) の exactness により, \\
(2.11) の列の適当な部分列は, 各項が有限群である限り
$2$-成分について exact になる.\\
% --- CHAPTER IV p.36 ---
列 \eqref{4.4} の部分列で,
\(\pit{n}{n},\ \pit{n+1}{n+1},\ \pit{2n+1}{2n+1},\ 
\pit{2n+1}{n},\ \pit{2n-1}{n}\)
が現れない箇所は, \\
\eqref{4.1} と,
$\Ex:\pis{n}{S^n}\xrightarrow{\cong}\pis{n+1}{S^{n+1}}$ が同型である事実から,
\[
\cdots\to\pit{n}{n}\xrightarrow{E}\pit{n+1}{n+1}\to\cdots
\]
の完全性を得る. 
残る問題は,
\[
\pit{2n}{n}\ \xrightarrow{\ \Ex\ }\ \pit{2n+1}{n+1}\
\xrightarrow{\ H\ }\ \pit{2n+1}{2n+1}\
\xrightarrow{\ \Delta\ }\ \pit{2n-1}{n}\
\xrightarrow{E} \ \pit{2n}{n+1}
\]
の exactness だが, これはLem\ref{lem:4.1} の証明中で本質的に示してある.\\ 
したがって, \eqref{4.4} 全体が exact である. 
\qed

\switchcolumn

\noindent
前ページ:\ $n$ を偶数とし, (2.11) の完全列を考える. \\
\(
\pis{2n-2}{S^{n-1}} \xrightarrow{\ \Ex\ } \pis{2n-1}{S^{n}}
\xrightarrow{\ H\ } \pis{2n-1}{S^{2n-1}}
\xrightarrow{\ \Delta\ } \pis{2n-3}{S^{n-1}}.
\)\\
$\Delta$ の象は $[\,\iota_{n-1},\iota_{n-1}\,]$ によって生成され, 
また, $2[\,\iota_{n-1},\iota_{n-1}\,]=0$ であるから, \\
\(
\pis{2n-1}{S^{n}} \cong
\mathbb{Z}\gen{\beta}\ \oplus\ 
E\pis{2n-2}{S^{n-1}} .
\)\\
ここで,
\(E(\pis{2n-2}{n-1})\) は有限であり, \\
$[\,\iota_{n-1},\iota_{n-1}\,]=0$ の場合, $H(\beta)=\iota_{2n-1}$,\\
$[\,\iota_{n-1},\iota_{n-1}\,]\ne0$ の場合, $H(\beta)=2\iota_{2n-1}$ である. 

\noindent
\textbf{Prop2.7.}\
$n$ が偶数のとき,
\(
H\!\bigl(\Delta(\iota_{2n+1})\bigr)\ =\ \pm\,2\,\iota_{2n-1}
\qquad (\text{cf.} (11.14)).
\)

\noindent
\textbf{Lem4.1.}\
$n$ が奇数のとき,  $\pi_{2n-1}^n=\pis{2n-1}{S^n;2} .$\\
$n$ が偶数で, $[\,\iota_{n-1},\iota_{n-1}\,]\ne0$ のとき, 
\(
\pi_{2n-1}^n\ \cong\ \mathbb{Z}\gen{\Delta(\iota_{2n+1})}\ \oplus \ 
\pis{2n-1}{S^n;2} .
\)\\
$n$ が偶数で, $[\,\iota_{n-1},\iota_{n-1}\,]=0$ のとき, 
\(
\pi_{2n-1}^n\ \cong\ \mathbb{Z}\gen{\alpha}\ \oplus \
\pis{2n-1}{S^n;2} .
\)\\
ここで, \(\alpha\) は 
\(H(\alpha)=\iota_{2n-1},\ E\alpha\in\pi_{2n}^{n+1}=\pi_{2n}(S^{n+1};2)\)
を満たす元である. \\
いずれの場合も, $[\,\iota_n,\iota_n\,]\in \pi_{2n-1}^n$ である. 

\noindent
\textbf{Th2.4.}\
$p$ を素数, $m>1$ とする. 
連続写像
\(
h:\ ((S^m)_\infty,(S^m)_{p-1})\longrightarrow ((S^{pm})_\infty,e_0)
\)\\
であって, $e^{pm}$ を $S^{pm}-e_0$ へ同相に移すものをとる. \\
i).\ $m$ が奇数で $p=2$ のとき, \\
\(
h_*:\ \pi_i\bigl((S^m)_\infty,S^m\bigr)\ \longrightarrow\ \pi_i\bigl((S^{2m})_\infty\bigr)
\)
は, 全ての $i$ に対して同型である. \\
ii).\ $m$ が偶数のとき, \\
\(
h_*:\ \pi_i\bigl((S^m)_\infty,(S^m)_{p-1}\bigr)\ \longrightarrow\ 
\pi_i\bigl((S^{pm})_\infty\bigr)
\)
は, $i<(p+1)m-1$ に対して同型であり, \\
さらに, 全ての $i$ に対して $p$-成分についての同型を与える. 

\noindent
\textbf{(2.11).}\ 
\(\cdots\ \longrightarrow\ 
\pi_i(S^m)\xrightarrow{\,E\,}\pi_{i+1}(S^{m+1})\xrightarrow{\,H\,}
\pi_{i+1}(S^{2m+1})\xrightarrow{\,\Delta\,}\pi_{i-1}(S^m)\xrightarrow{\,E\,}\cdots
\)\\
これは $m$ が奇数, あるいは $i<3m-1$ のとき完全列である. \\
$m$ が偶数で, $i\ge 3m-1$ の場合には, (2.11) を
2-成分に関する完全列として考える. \\
\(\pis{2n+1}{S^{2n+1}}\xlongrightarrow{\Delta}
\pis{2n-1}{S^{n}}\xlongrightarrow{E}
\pis{2n}{S^{n+1}}\).\\
\(\ker E = \im \Delta = \gen{\Delta(\iota_{2n+1})}
\subset\pis{2n-1}{S^n; 2}\).


\end{paracol}

\begin{paracol}{2}

素数 $p$ に対し, 安定群 $G_k$ の $p$-成分を $(G_k;p)$ と書く. \\
\eqref{4.4} の exactness と (3.2) より, 次の同型を得る：
\begin{equation}\label{4.5}
% (4.5)
\Ex^{\,m-n}:\pit{n+k}{n}\xrightarrow{\cong}\pit{m+k}{m}\quad(m\ge n),\ \ \
\Ex^{\infty}:\pit{n+k}{n}\xrightarrow{\cong}(G_k;2)
\quad(n\ge k+2).
\end{equation}

\begin{lemma}\label{lem:4.3}
% Lem4.3
$n\ge1$ とする. $\alpha,\gamma$ が $p$-成分に属するならば, \\
$\alpha\circ \Ex^{n}\beta$ と $\beta\circ\gamma$ も, 
$p$-成分に属する. \\
さらに, $\alpha\circ \Ex^{n}\beta=\beta\circ\gamma=0$ のとき, 
$\todai[n]{\alpha,\Ex^{n}\beta,\Ex^{n}\gamma}$ の各元は,
$p$-成分に属する. 
\end{lemma}

この Lemma は, 球面のホモトピー群に対して述べたが, 一般の空間に対しても成り立つ：
\begin{equation}\label{4.6}
% (4.6).\
\raisetag{3ex}
\begin{aligned}[t]
  &n\ge 1,\ p^r\alpha=0,\ p^s\gamma=0,\ r,s>0,\ 
  \alpha\in[E^n Y, Z],\ \gamma\in[EW, X] \ \text{のとき, }\\
  &\beta\in[X, Y]\ \text{に対して, } \ 
  p^r(\alpha\circ E^r\beta) = 0,\ p^s(\beta\circ\gamma) = 0 \ \text{である. }\\
  &\text{さらに, }\ \alpha\circ E^n\beta=0,\ \beta\circ\gamma = 0 \ \text{のとき, }
  p^{r+s}\todai[n]{\alpha, E^n\beta, E^n\gamma} = 0 \ \text{である. }\qquad
\end{aligned}
\end{equation}
%
\begin{equation}\label{4.7}
% (4.7).\
\raisetag{3ex}
\begin{aligned}[t]
  & n\ge1, \ \alpha\in\pis{k+n}{S^m; 2},\ \gamma\in\pis{i}{S^j; 2}
    \ \text{のとき, }\\
  & \alpha\circ E^n\pis{j}{S^k} = \alpha\circ E^n \pit{j}{k}
  \subset \pit{j+n}{m},
  \ \pis{j}{S^k}\circ\gamma = \pit{j}{k}\circ\gamma \subset\pit{i}{k}\ . \\
  &\text{さらに, }\ \alpha\circ E^n\beta=0,\ \beta\circ\gamma = 0 \ \text{のとき, }\\
  & \todai[n]{\alpha, E^n\beta, E^n\gamma}\ \text{は,}\
  \alpha\circ E^n\pit{i+1}{k} + \pit{j+n+1}{m}\circ E^{n+1}\gamma\
  \text{の coset で, }\qquad\qquad\qquad\\
  & \pit{i+n+1}{m} \ \text{の部分集合である. }
\end{aligned}
\end{equation}

% (p=2 の特別化 4.7 は次頁で展開) 

% --- CHAPTER IV p.37-38 ---

\switchcolumn

\noindent
\textbf{(3.2).}\
\(\Ex: \pis{i}{S^m} \longrightarrow \pis{i+1}{S^{m+1}}
\quad\text{は, } i<2m-1 \text{で同型, } i=2m-1 \text{で全射である. }\)



\end{paracol}

\begin{paracol}{2}

\noindent
$n=2,4,8$ のとき, Hopf fibration
$S^n\hookrightarrow S^{2n-1}\xrightarrow{h} S^n$ が存在し, これに付随して,
\[
E+h_*:\ \pis{i-1}{S^{n-1}}\oplus\pis{i}{S^{2n-1}}
\ \xrightarrow{\ \cong\ }\ \pis{i}{S^n}
\]
が, 任意の \(i\) で同型になることが知られている \cite{Steenrod1951}.\\
これは, $h$ を Hopf 不変量 $\pm1$ を持つ写像に置き換えても成り立つ. \\
なぜなら, これらの同型は, 次の写像
\[
i\cdot\Omega h: S^{n-1}\times\Ox S^{2n-1}\to\Ox S^n
\]
によって誘導されるからである. 
ここで, "\(\cdot\)" は \(\Omega S^n\) の道の積を表す. \\
この写像は, コホモロジーの同型を誘導するので,ホモロジーとホモトピーの同型を誘導する. \\
証明の詳細は読者に任せる. 

\begin{proposition}\label{prop:4.4}
% Prop4.4.\
$\alpha\in\pit{2n-1}{n}$ で $H(\alpha)=\pm\iota_{2n-1}$ とする. 
このとき,\\
\(
\pit{i-1}{n-1}\oplus\pit{i}{2n-1}\to\pit{i}{n},\quad
(\beta,\gamma) \longmapsto E\beta + \alpha\circ\gamma
\)
は, 任意の \(i\) で同型である. \\
特に, 
\(
E:\pit{i-1}{n-1}\to \pit{i}{n}
\)
は, 任意の \(i\) で単射である. 
\end{proposition}

\switchcolumn

\noindent
\textbf{
\cite{Steenrod1951}
N.~E.~Steenrod,}
\newblock Topology of fibre bundles,
\newblock Princeton.

\end{paracol}

\begin{paracol}{2}

\begin{lemma}\label{lem:4.5}
% Lem4.5
$n=2,4,8$ とする. \\
\(
E:\pis{i-1}{S^{n-1}}\to\pis{i}{S^n}
\)
は, 任意の $i$ で単射であり, \\
任意の整数 $r$ と $\beta\in\pis{i-1}{S^{n-1}}$ に対して, 
\[
r\,\iota_{n-1}\circ \beta = r\beta \ .
\]
\end{lemma}

\noindent
\textbf{Proof.}\
前半は, \(\pis{i}{S^n}\) の Hopf 分解より明らかである. \\
(2.1) により, $E(r\,\iota_{n-1}\circ\beta)=E(r\beta)$ なので, 後半も成り立つ. 
\qed

\medskip
\noindent
Lem4.1 の証明から, 次が分かる. 
\begin{equation}\label{4.8}
% (4.8).\
\raisetag{3ex}
\begin{aligned}[t]
& n \ \text{を偶数とする. } \\
& \Delta(\iota_{2n-1})=[\iota_{n-1},\iota_{n-1}]\ne0 \ \text{のとき, }\\
& \pit{2n-1}{n}
=E\,\pit{2n-2}{\,n-1}\ \oplus\ \Delta\,\pit{2n+1}{\,2n+1}
\cong E\,\pit{2n-2}{\,n-1}\ \oplus\ \mathbb{Z},\ \pit{2n}{\,n+1} = E^{2}\,\pit{2n-2}{\,n-1} 
\cong E\,\pit{2n-2}{\,n}.\\
& \Delta(\iota_{2n-1})=0 \ \text{のとき, }\\
& E^{2}:\pit{2n-2}{\,n-1}\to\pit{2n}{\,n+1}\ \text{は単射で, }\
\pit{2n}{\,n+1}/E^{2}\pit{2n-2}{\,n-1}\cong\mathbb{Z}_2 . \\
& \text{さらに,}\ \alpha\in\pit{2n-1}{n} \ \text{に対して, 次が成り立つ. }\\
& E\alpha\in E^{2}\pit{2n-2}{\,n-1} \ \Leftrightarrow \
H(\alpha)\in 2\,\pit{2n-1}{\,2n-1}\ .
\end{aligned}
\end{equation}

\begin{equation}\label{4.9}
% \noindent
% \textbf{(4.9).}\
\pit{m}{n}\circ\pit{k}{m}\ \subset\ \pit{k}{n}.
\end{equation}

\noindent
\textbf{Proof.}\
$k\ne m,\,2m-1$ のとき, $\pit{k}{m}=\pi_k(S^{m};2)$ なので, 
\eqref{4.7} から, (4.9) が成り立つ. 

$k=m$ のとき,  $\pit{m}{n}\circ\pit{m}{m} = \pit{m}{n}$ なので, 
(4.9) が成り立つ. 

$k=2m-1,\ m$ が奇数のとき, 
Lem \ref{lem:4.1} より, 
\(\pit{k}{m} = \pis{k}{S^m; 2}\)
なので, \\
(4.7) から, (4.9) が成り立つ. 

$k=2m-1,\ m$ が偶数のとき, 
$\alpha\in\pit{2m-1}{m}$ を考える. \\
$H[\iota_m,\iota_m] = \pm 2\,\iota_{2m-1}$ なので, ある整数 $r$ に対し,
$H(2\alpha)=H(r[\iota_m,\iota_m])$ である. \\
(4.4) の完全性から, 
\(2\alpha=r[\iota_m,\iota_m]+E\gamma\)
となる $\gamma\in\pit{2m-2}{\,m-1}$ が存在する. \\
任意の $\beta\in\pit{m}{n}$ をとり, 
$\beta$ の位数を $2^{t}$ とする. \\
このとき, \(2^{2t+1}(\beta\circ\alpha) = \) を示す. 
(2.1) と, Whitehead 積の双線型性と自然性 \cite{Whitehead1941} から, 
\[
\begin{aligned}
2^{2t+1}(\beta\circ\alpha)
&=\beta\circ r 2^{2t}[\iota_m,\iota_m]
\ +\ \beta\circ 2^{2t}E\gamma\\
&=\beta\circ r\bigl[2^{t}\iota_m,\,2^{t}\iota_m\bigr]\ +\ 2^{2t}\beta\circ E\gamma\\
&=\beta\circ 2^{t}\iota_m\circ r[\iota_m,\iota_m]\ +\ 0\\
&= 2^{t}\beta\circ r[\iota_m,\iota_m]\;=\;0 .
\end{aligned}
\]
よって, $\beta\circ\alpha$ は $2$-成分の元であり, 
$\pit{m}{n}\circ\pit{2m-1}{m}\subset\pit{2m-1}{n}$ が成り立つ. 
\qed

\switchcolumn

\noindent
\textbf{(4.7).}\
\(n\ge1, \ \alpha\in\pis{k+n}{S^m; 2},\ \gamma\in\pis{i}{S^j; 2}
    \ \text{のとき, }\)\\
\(\alpha\circ E^n\pis{j}{S^k} = \alpha\circ E^n \pit{j}{k}
  \subset \pit{j+n}{m},
  \ \pis{j}{S^k}\circ\gamma = \pit{j}{k}\circ\gamma \subset\pit{i}{k}\ . \)\\
\(\text{さらに, }\ \alpha\circ E^n\beta=0,\ \beta\circ\gamma = 0 \ \text{のとき, }\)\\
\(\todai[n]{\alpha, E^n\beta, E^n\gamma}\ \text{は,}\
  \alpha\circ E^n\pit{i+1}{k} + \pit{j+n+1}{m}\circ E^{n+1}\gamma\
  \text{の coset で, }\pit{i+n+1}{m} \ \text{の部分集合である. }\)

\noindent
\textbf{Lem4.1.}\
$n$ が奇数のとき,  $\pi_{2n-1}^n=\pis{2n-1}{S^n;2} .$\\
$n$ が偶数で, $[\,\iota_{n-1},\iota_{n-1}\,]\ne0$ のとき, 
\(
\pi_{2n-1}^n\ \cong\ \mathbb{Z}\gen{\Delta(\iota_{2n+1})}\ \oplus \ 
\pis{2n-1}{S^n;2} .
\)\\
$n$ が偶数で, $[\,\iota_{n-1},\iota_{n-1}\,]=0$ のとき, 
\(
\pi_{2n-1}^n\ \cong\ \mathbb{Z}\gen{\alpha}\ \oplus \
\pis{2n-1}{S^n;2} .
\)\\
ここで, \(\alpha\) は 
\(H(\alpha)=\iota_{2n-1},\ E\alpha\in\pi_{2n}^{n+1}=\pi_{2n}(S^{n+1};2)\)
を満たす元である. \\
いずれの場合も, $[\,\iota_n,\iota_n\,]\in \pi_{2n-1}^n$ である. 

\noindent
\textbf{(4.4).}\
\(\pit{2m-2}{m-1}\xlongrightarrow{E}
\pit{2m-1}{m}\xlongrightarrow{H}
\pit{2m-1}{2m-1}\)\\
\qquad\qquad\(\gamma\longmapsto 2\alpha-[\iota_m, \iota_m]\longmapsto 0\)

\noindent
\textbf{\cite{Whitehead1941}
J.~H.~C.~Whitehead,}
\newblock On adding relations to homotopy groups,
\newblock Ann. of Math., {\bf 42} (1941), 409--428.

\noindent
\([\alpha_1+\alpha_2, \beta] =
[\alpha_1, \beta] + [\alpha_2, \beta],\)\\
\([\alpha, \beta_1 + \beta_2] =
[\alpha, \beta_1] + [\alpha, \beta_2],\)\\
\([\gamma\circ\alpha, \gamma\circ\beta] = 
\gamma\circ[\alpha, \beta].\)



\end{paracol}

\newpage

\begin{paracol}{2}

\renewcommand{\refname}{References}
\begin{thebibliography}{99}

\bibitem{Adams1960}
J.~F.~Adams,
\newblock On the non-existence of elements of Hopf invariant one,
\newblock Ann. of Math., {\bf 72} (1960), 20--104.

\bibitem{Adem}
J.~Adem,
\newblock The relations on Steenrod powers of cohomology classes,
\newblock Algebraic geometry and topology, Princeton.

\bibitem{Barratt1955}
M.~G.~Barratt,
\newblock Track groups I, II,
\newblock Proc. London Math. Soc., {\bf 5} (1955), 71--106, 285--329.

\bibitem{Barratt1961}
M.~G.~Barratt,
\newblock Note on a formula due to Toda,
\newblock Jour. London Math. Soc., {\bf 36} (1961), 95--96.

% 5
\bibitem{BarrattHilton1953}
M.~G.~Barratt and P.~J.~Hilton,
\newblock On join operations of homotopy groups,
\newblock Proc. London Math. Soc., {\bf 3} (1953), 430--445.

\bibitem{BlakersMassey1952}
A.~L.~Blakers and W.~S.~Massey,
\newblock The homotopy groups of a triad II,
\newblock Ann. of Math., {\bf 55} (1952), 192--201.

\bibitem{Freudenthal1937}
H.~Freudenthal,
\newblock Über die Klassen der Sphärenabbildungen,
\newblock Compos. Math., {\bf 5} (1937).

\bibitem{Hopf1935}
H.~Hopf,
\newblock Über die Abbildungen von Sphären auf Sphären niedriger Dimension,
\newblock Fund. Math., {\bf 25} (1935).

\bibitem{James1955}
I.~M.~James,
\newblock Reduced product spaces,
\newblock Ann. of Math., {\bf 62} (1955), 170--197.

% 10
\bibitem{James1956}
I.~M.~James,
\newblock Suspension triad of a sphere,
\newblock Ann. of Math., {\bf 63} (1956), 407--429.

\bibitem{James1958}
I.~M.~James,
\newblock The intrinsic join: A study of the homotopy groups of Stiefel manifold,
\newblock Proc. London Math. Soc., {\bf 8} (1958), 507--535.

\bibitem{Moore1953}
J.~C.~Moore,
\newblock Some applications of homology theory to homotopy problems,
\newblock Ann. of Math., {\bf 58} (1953), 325--350.

\bibitem{Serre1951}
J.-P.~Serre,
\newblock Homologie singulière des espaces fibrés,
\newblock Ann. of Math., {\bf 54} (1951), 425--505.

\bibitem{Serre1953a}
J.-P.~Serre,
\newblock Cohomologie modulo 2 des complexes d'Eilenberg-MacLane,
\newblock Comm. Math. Helv., {\bf 27} (1953), 198--231.

% 15
\bibitem{Serre1953b}
J.-P.~Serre,
\newblock Groups d'homotopie et classes groupes abéliens,
\newblock Ann. of Math., {\bf 58} (1953), 258--294.

\bibitem{Serre1953c}
J.-P.~Serre,
\newblock Quelques calculs de groupes d'homotopie,
\newblock Comptes Rendus Paris, {\bf 236} (1953), 2475--2477.

\bibitem{Steenrod1951}
N.~E.~Steenrod,
\newblock Topology of fibre bundles,
\newblock Princeton.

\bibitem{Steenrod1949}
N.~E.~Steenrod,
\newblock Cohomology invariants of mappings,
\newblock Ann. of Math., {\bf 50} (1949), 954--988.

\bibitem{Steenrod1953}
N.~E.~Steenrod,
\newblock Cyclic reduced powers of cohomology classes,
\newblock Proc. Nat. Acad. Sci. U.S.A., {\bf 39} (1953), 217--223.

% 20
\bibitem{Toda1952}
H.~Toda,
\newblock Generalized Whitehead products and homotopy groups of spheres,
\newblock Jour. of Inst. Poly. Osaka City Univ., {\bf 3} (1952), 43--82.

\bibitem{Toda1956}
H.~Toda,
\newblock On double suspension $E^2$,
\newblock Jour. of Inst. Poly. Osaka City Univ., {\bf 7} (1956), 103--145.

\bibitem{Toda1957}
H.~Toda,
\newblock Reduced join and Whitehead product,
\newblock Jour. of Inst. Poly. Osaka City Univ., {\bf 8} (1957), 15--30.

\bibitem{Toda1958a}
H.~Toda,
\newblock On exact sequences in Steenrod algebra mod 2,
\newblock Memoirs Univ. of Kyoto, {\bf 31} (1958), 33--64.

\bibitem{Toda1958b}
H.~Toda,
\newblock $p$-primary components of homotopy groups II. mod $p$ Hopf invariant,
\newblock Memoirs Univ. of Kyoto, {\bf 31} (1958), 143--160.

% 25
\bibitem{Toda1959}
H.~Toda,
\newblock $p$-primary components of homotopy groups IV. Compositions and toric constructions,
\newblock Memoirs Univ. of Kyoto, {\bf 32} (1959), 297--332.

\bibitem{Whitehead1950a}
G.~W.~Whitehead,
\newblock A generalization of the Hopf invariant,
\newblock Ann. of Math., {\bf 51} (1950), 192--237.

\bibitem{Whitehead1950b}
G.~W.~Whitehead,
\newblock The $(n+2)$nd homotopy group of the $n$-sphere,
\newblock Ann. of Math., {\bf 52} (1950), 245--247.

\bibitem{Whitehead1941}
J.~H.~C.~Whitehead,
\newblock On adding relations to homotopy groups,
\newblock Ann. of Math., {\bf 42} (1941), 409--428.

\bibitem{Whitehead1944}
J.~H.~C.~Whitehead,
\newblock On the groups $\pi_n(V_{m,n})$ and sphere bundles,
\newblock Proc. London Math. Soc., {\bf 48} (1944), 243--291.

\end{thebibliography}


\end{paracol}

\end{document}
